\chapter{Arnold's topological proof}
\label{chap:arnold}

This proof is not formalized yet. The main theorem is stated in
\texttt{Quintic/Arnold.lean} with a \texttt{sorry}. The lemmas below are the planned
decomposition; none has a Lean counterpart yet. Each one, once stated in Lean, should get a
\texttt{\textbackslash lean\{\dots\}} tag here, and \texttt{\textbackslash leanok} once proved.

The argument works over the space of complex coefficient vectors
$c = (c_0, \dots, c_4) \in \mathbb{C}^5$ and the polynomial
$P_c(X) = X^5 + \sum_i c_i X^i$.

\section{Monodromy of the roots}

\begin{definition}[Polynomials with distinct roots]
  \label{def:config-space}
  Let $U \subseteq \mathbb{C}^5$ be the set of $c$ such that $P_c$ has five distinct roots,
  i.e.\ the complement of the discriminant hypersurface $\{\Delta = 0\}$.
\end{definition}

\begin{lemma}[Roots form a covering]
  \label{lem:roots-covering}
  \uses{def:config-space}
  Over $U$, the set of roots $\{(c, z) : P_c(z) = 0\}$ is a $5$-sheeted covering space of $U$,
  and $U$ is path-connected.
\end{lemma}

\begin{lemma}[Monodromy of the roots is $S_5$]
  \label{lem:roots-monodromy}
  \uses{lem:roots-covering}
  Lifting loops in $U$ based at $c_0$ permutes the five roots of $P_{c_0}$, and every
  permutation arises this way.
\end{lemma}

\begin{proof}
  Choose the roots first: moving two roots around each other while fixing the others gives
  a transposition, and the coefficients follow by Vieta. Transpositions generate $S_5$.
\end{proof}

\section{Monodromy of radical expressions}

\begin{definition}[Branches of a formula]
  \label{def:branches}
  \uses{def:radexpr}
  For a formula $e$, let $Z_e \subseteq \mathbb{C}^5$ be the proper analytic subset where some
  subformula that is inverted or has a root taken vanishes. On $\mathbb{C}^5 \setminus Z_e$, the
  values of $e$ form a finite-sheeted covering (its branches).
\end{definition}

\begin{lemma}[Complements of proper analytic subsets]
  \label{lem:complement-connected}
  The complement of a proper analytic subset of a connected complex manifold is connected,
  and every loop in $U$ is homotopic to one avoiding any given proper analytic subset.
\end{lemma}

\begin{lemma}[Monodromy of a formula is solvable]
  \label{lem:radical-monodromy-solvable}
  \uses{def:branches, lem:complement-connected}
  The monodromy group of the branches of any formula in radicals is solvable.
\end{lemma}

\begin{proof}
  Induction on the formula. Sums, products and inverses act on pairs of branches, so the
  group embeds in a product of solvable groups. For $\sqrt[n]{f}$, the kernel of the map to the
  monodromy group of $f$ multiplies branches by $n$-th roots of unity, so it is abelian, and an
  extension of a solvable group by an abelian group is solvable.
\end{proof}

\begin{lemma}[$A_5$ is perfect]
  \label{lem:a5-perfect}
  The alternating group $A_5$ equals its commutator subgroup. Consequently $S_5$ is not
  solvable, and every iterated commutator subgroup of $S_5$ contains $A_5$.
\end{lemma}

\section{The bridge}

\begin{lemma}[From algebraic values to analytic branches]
  \label{lem:bridge}
  \uses{def:value, def:branches}
  Fix $c_0$ outside $Z_e$ and the discriminant locus. If $\mathrm{Value}(e, y)$ with $y \in L$,
  there is a branch of $e$ near $c_0$, a germ of a holomorphic function $g$, such that any
  polynomial identity over $K$ satisfied by $y$ is satisfied by $g$. In particular, if $P(y) = 0$
  then $P_c(g(c)) = 0$ near $c_0$.
\end{lemma}

\begin{proof}
  Embed the radical tower generated by the subformula values of $e$ into the field of
  meromorphic germs at $c_0$, choosing a branch of each root compatibly.
  Alternatively, restate \texttt{QuinticUnsolvable} analytically and move this bridge into the
  Galois proof.
\end{proof}

\section{Conclusion}

\begin{theorem}[Abel--Ruffini, Arnold's proof]
  \label{thm:arnold-main}
  \lean{Quintic.quintic_unsolvable_arnold}
  \uses{def:quintic-unsolvable, lem:roots-monodromy, lem:radical-monodromy-solvable,
    lem:a5-perfect, lem:bridge}
  The quintic is unsolvable.
\end{theorem}

\begin{proof}
  Suppose $\mathrm{Value}(e, y)$ and $P(y) = 0$. By \cref{lem:bridge}, a branch $g$ of $e$ is a
  root of $P_c$ near $c_0$. By \cref{lem:radical-monodromy-solvable} some iterated commutator of
  loops fixes every branch of $e$, while by \cref{lem:roots-monodromy} and \cref{lem:a5-perfect}
  the same iterated commutators of loops act on the roots through $A_5$, so one of them moves
  the root $g(c_0)$. Continuing $g$ along that loop both fixes it and moves it, a contradiction.
\end{proof}
